\subsubsection{An exponentially small central heat-tail bound}
\label{subsec:sixband-exponential}
The six-band comparison in Theorem~\ref{thm:sixband}
gives a nonnegative defect that is supported, under its
explicit upper envelope, only outside the interval
$[-6,6]$. The polynomial bound used in
Corollary~\ref{cor:sixbandduhamel} discards
the Gaussian concentration of the heat semigroup.
Retaining it yields a substantially sharper,
but still nonzero, source estimate.

\begin{theorem}[Exponential six-band Duhamel bound]
\label{thm:sixband-exponential}
Retain $T=1/2$, $k_n=(n+1)(2n+1)$ and
$A_{n+1}(t)=L_{n+1,t}(0)$, as in
Theorem~\ref{thm:sixband}. Let
\[
E_n(s,x)=
[k_nU_{n+1}(s,x)-c_n(s)U_n(s,x)]_+,
\quad
c_n(s)=4k_n\frac{A_{n+1}(T-s)}{A_n(T-s)}.
\]
Then every $n\ge1$ satisfies
\begin{equation}\label{eq:sixband-exponential}
\boxed{
\begin{aligned}
0\le
\int_0^T(P_{T-s}E_n(s,\cdot))(0)\,ds
&\le \frac8{45}e^{-135/4}
\,k_n\sup_{0\le t\le T}A_{n+1}(t)\\
&<
\frac{k_n}{10^{14}}
\sup_{0\le t\le T}A_{n+1}(t).
\end{aligned}}
\end{equation}
The right side improves the coefficient $1/3456$
of Corollary~\ref{cor:sixbandduhamel} by
more than ten orders of magnitude.
\end{theorem}
\begin{proof}
Put $b(x)=(x^2/12-3)_+=(x^2-36)_+/12$.
For $a=15/16$ and $v\ge0$,
$av\le e^{av}$, so for all real $x$,
\begin{equation}\label{eq:sixband-pointwise}
b(x)\le
\frac1{12a}e^{a(x^2-36)}
=\frac4{45}e^{-135/4}e^{15x^2/16}.
\end{equation}
Let $Z_h$ be a centred real Gaussian
with variance $h\in[0,1/2]$.
Completing the square in the Gaussian
integral gives
\[
\mathbb E e^{aZ_h^2}=(1-2ah)^{-1/2}.
\]
Since $1-2ah\ge1/16$, equation
\eqref{eq:sixband-pointwise} implies
\begin{equation}
(P_hb)(0)\le
\frac{16}{45}e^{-135/4}
\qquad(0\le h\le1/2).
\end{equation}
The six-band source inequality states
\[
E_n(s,x)\le k_nA_{n+1}(T-s)b(x).
\]
Positivity of $P_{T-s}$, integration over
$0\le s\le T$ and the preceding Gaussian
bound yield the first line of
\eqref{eq:sixband-exponential}.

For the rational final comparison,
$135/4>33$ and the positive Taylor series
gives $e^3>20$. Consequently
\[
\frac8{45}e^{-135/4}
<
\frac8{45\cdot20^{11}}
=\frac1{1\,152\,000\,000\,000\,000}
<10^{-14}.
\]
No floating-point comparison enters
this part of the proof.
\end{proof}

\paragraph{A spatial envelope and its limitation.}
For an observation point $x_0\in\mathbb R$,
the noncentral Gaussian identity
\[
\mathbb E e^{a(x_0+Z_h)^2}
=(1-2ah)^{-1/2}
\exp\!\left(\frac{a x_0^2}{1-2ah}\right)
\]
gives the further estimate
\begin{equation}\label{eq:sixband-spatial-upper}
\int_0^{1/2}
(P_{1/2-s}E_n(s,\cdot))(x_0)\,ds
\le \frac8{45}
e^{-135/4+15x_0^2}\,
k_n\sup_t A_{n+1}(t).
\end{equation}
This is useful near the central observation point,
but it grows in $|x_0|$ and does not
establish the global sign of $U_n$.

Crucially, even a very small positive
upper bound on the heat-smoothed
comparison defect does \emph{not}
show that the defect vanishes. The
RH-equivalent criterion of
Theorem~\ref{thm:onepointcriterion}
demands a uniform bound under an
unbounded penalty, or an exact
zero-defect statement, rather than
a fixed small nonzero upper bound.
Moreover the factor
$k_n\sup_t A_{n+1}(t)$ is
order-dependent, and no
all-order relative sign conclusion
follows from \eqref{eq:sixband-exponential}.
The result is a quantitative
improvement of the completed
six-band estimate, not a proof of RH.

The rational identities and
elementary exponential comparisons
are reproduced by
\path{certificates/xi_six_band_exponential_exact.py};
the corresponding Lean algebraic
gates are in
\path{XiKernel/SixBandExponentialTail.lean}.
The Gaussian moment formula
and the genuine theta-correlation
source inequality are analytic
inputs to this argument, not
end-to-end Lean formalizations.
